% Данный файл распространяется под лицензией CC BY 4.0. Текст лицензии размещён на https://creativecommons.org/licenses/by/4.0/
% (c) Кафедра системного программирования, 2026

% Компилировать XeTeX

\documentclass[a4paper]{article}

\usepackage[a4paper, top=8mm, bottom=8mm, left=8mm, right=8mm]{geometry}

\usepackage{polyglossia}
\setdefaultlanguage[babelshorthands=true]{russian}
\setotherlanguage{english}

\usepackage{fontspec}
%\setmainfont{FreeSerif}
\setmainfont{DejaVu Serif}
\newfontfamily{\russianfonttt}[Scale=0.7]{DejaVuSansMono}

\usepackage[tiny, compact]{titlesec}

\usepackage{titling}
\setlength{\droptitle}{-1cm}
\pretitle{\begin{center}\begin{bfseries}\Large}
\posttitle{\par\end{bfseries}\end{center}}
\preauthor{\begin{center}\normalsize}
\postauthor{\par\end{center}\vspace{-1.8cm}}

\usepackage{hyperref}
\usepackage{bookmark}
\usepackage{csquotes}

\title{Архиватор на базе алгоритма Хаффмана}

\author{Анпилогов Дарий Михайлович}

\date{}

\begin{document}

\maketitle

\begin{flushright}
    Группа: \emph{25.Б41-мм}

    Кафедра: \emph{системного программирования}

    Научный руководитель: \emph{Литвинов Юрий Викторович}

    Номер семестра практики: \emph{2}
\end{flushright}

\section{Введение}

Написание архиватора, основанного на базе алгоритма Хаффмана, "--- это учебная задача, направленная на получение опыта написания курсовой работы, изучения стандартных требований, процессов, а также способов достижения поставленных в рамках работы целей.
Несмотря на наличие подобных решений, понимание проблем и требований при их разработке является важной частью обучения.

\section{Постановка задачи}

\textbf{Целью работы является} реализация консольного приложения, выполняющего в зависимости от параметров командной строки сжатие или разжатие файла по алгоритму Хаффмана.

\textbf{Для достижения цели требуется решить следующие задачи:}
\begin{enumerate}
    \item Выполнить обзор существующих решений.
    \item Выбрать инструменты для реализации.
    \item Спроектировать решение.
    \item Реализовать решение.
    \item Выполнить экспериментальное исследование.
\end{enumerate}

\section{Обзор}

Целью обзора является изучение алгоритма Хаффмана и существующих подходов к его реализации для выбора архитектурных решений разрабатываемого архиватора.

Алгоритм Хаффмана был предложен Дэвидом Хаффманом в 1952 году и предназначен для сжатия данных без потерь. Алгоритм строит префиксный код переменной длины на основе частот появления символов во входных данных. Более часто встречающимся символам назначаются более короткие коды, а редким --- более длинные. Для заданного набора частот алгоритм обеспечивает минимальную среднюю длину кодового слова среди всех префиксных кодов.

Алгоритм Хаффмана применяется как самостоятельный метод сжатия, а также используется в составе более сложных алгоритмов и форматов хранения данных. В частности, кодирование Хаффмана входит в состав алгоритма DEFLATE, применяемого в форматах ZIP, gzip и PNG, а также используется на этапе энтропийного кодирования в форматах JPEG и MP3.

При реализации архиватора необходимо решить две основные задачи: построение дерева Хаффмана и сохранение информации, необходимой для восстановления кодов при распаковке. Для построения дерева в большинстве реализаций используется очередь с приоритетами, позволяющая эффективно выбирать узлы с минимальными частотами. В данной работе используется контейнер \texttt{std::priority\_queue}, обеспечивающий построение дерева за $O(n \log n)$, где (n) --- количество различных символов входных данных.

Для восстановления кодов при распаковке распространены два подхода: сохранение таблицы частот символов либо сохранение непосредственно кодов. В данной работе выбран второй вариант, поскольку он позволяет восстановить таблицу кодов без хранения полной таблицы частот и уменьшает объём служебной информации в архиве.

Таким образом, алгоритм Хаффмана является подходящей основой для реализации учебного архиватора благодаря сравнительно простой реализации, использованию эффективных структур данных и широкому распространению в современных форматах хранения данных, а также позволяет изучить работу с двоичными данными, файловым вводом-выводом, тестированием программного обеспечения и процессом разработки законченного программного решения.

\section{Описание предлагаемого решения}

Для реализации проекта были выбраны:
\begin{itemize}
    \item Язык программирования: C++17. Причина "--- скорость и простота работы с потоками данных и битовыми операциями.
    \item Система сборки: CMake v3.14. Причина "--- простота сборки, кроссплатформенность.
    \item Тестирование: GoogleTest v1.14.0. Причина "--- стандарт индустрии, функциональность.
    \item CI: GitHub Actions. Причина "--- простота интеграции с платформой распространения проекта. Используется для запуска тестов и проверки программы с использованием clang-format, AddressSanitizer и UndefinedBehaviorSanitizer.
\end{itemize}

Реализация разделена на 2 части:
\begin{itemize}
    \item Библиотека, реализующая функции архивации и разархивации.
    \item Консольное приложение, использующее библиотеку и предоставляющее пользователю интерфейс для использования функций.
\end{itemize}

Основными компонентами библиотеки являются:
\begin{itemize}
    \item \texttt{BitReader} "--- чтение отдельных битов из входного потока с использованием внутреннего буфера;
    \item \texttt{BitWriter} "--- запись отдельных битов в выходной поток;
    \item \texttt{Bits} "--- контейнер для хранения битовых последовательностей;
    \item \texttt{CodesTable} "--- таблица кодов Хаффмана для всех возможных значений байта;
    \item \texttt{CountTable} "--- таблица частот встречаемости байтов (количества вхождений каждого байта);
    \item \texttt{archive()} "--- реализация алгоритма сжатия;
    \item \texttt{unarchive()} "--- реализация алгоритма разжатия.
\end{itemize}

Для построения кодов Хаффмана используется очередь с приоритетами \texttt{std::priority\_queue}. Дерево в явном виде не строится. На каждой итерации алгоритма склеиваются два набора байтов с минимальными суммарными вхождениями в файл. К кодам соответствующих байтов добавляется по одному биту в зависимости от множества, к которому они принадлежат.
Для разархивации явным образом строится дерево по таблице частот.

Архив содержит размер таблицы кодов, описание таблицы кодов, размер сжатых данных и сжатые данные. Не содержащиеся в исходном файле коды не хранятся в таблице кодов. Такой подход позволяет восстановить таблицу кодов без необходимости хранить полную таблицу частот.

Поддерживаются два режима работы архиватора. В режиме \texttt{twice} входной файл читается дважды: первый проход используется для подсчёта частот, второй "--- для кодирования данных. Данный режим позволяет обрабатывать файлы произвольного размера без полного размещения их содержимого в памяти. В режиме \texttt{ram} файл полностью загружается в оперативную память, что уменьшает количество операций ввода-вывода и позволяет работать с потоками, не поддерживающими повторное чтение.

Приложение поддерживает обработку ошибочных ситуаций: неверные аргументы командной строки, отсутствие входного файла, ошибки чтения и записи, а также повреждённые архивы. Для завершения работы используются различные коды возврата: 0 при успешном выполнении, 1 при ошибке параметров командной строки и -1 при прочих ошибках.

Для релизной версии дополнительно используются межпроцедурная оптимизация и оптимизации компилятора. Проект распространяется по лицензии Apache License 2.0.

Реализовано как модульное, так и интеграционное тестирование "--- всего 47 тестов.

\section{Эксперименты}

Для проведения экспериментального исследования были выбраны следующие инструменты:
\begin{itemize}
    \item Язык программирования: Python 3.13. Причина "--- простота кода, совместимость с различными типами данных.
    \item Виртуальное окружение: uv. Причина "--- возможность явно указать используемые библиотеки, простой запуск.
    \item Формат сохранения данных: CSV. Причина "--- простота чтения и записи данных.
    \item Библиотека отрисовки графиков: matplotlib. Причина "--- стандарт для отрисовки графиков, простота использования, распространённость.
\end{itemize}

Для исследования производительности был разработан набор автоматизированных экспериментов. Измерения проводились на компьютере с процессором Intel Core i7-9750H, 16 ГиБ оперативной памяти под управлением операционной системы openSUSE Tumbleweed.

В качестве тестовых данных использовались:
\begin{itemize}
    \item текстовые файлы;
    \item случайные данные;
    \item периодические последовательности с периодом 1, 2, 4, 8, 16 и 32 байта.
\end{itemize}

Размер файлов изменялся от 1 байта до 128 МиБ по степеням двойки. Для каждого измерения выполнялось 10 независимых запусков. Рассчитывались математическое ожидание и среднеквадратичное отклонение времени выполнения, а также скорости сжатия и разжатия.

Результаты показали, что степень сжатия существенно зависит от энтропии входных данных. Для периодических последовательностей коэффициент сжатия стремится к значению 0.125, что соответствует восьмикратному уменьшению размера файла. Для данных с высокой энтропией эффективность сжатия снижается, а в отдельных случаях размер архива может превышать размер исходного файла из-за накладных расходов на хранение таблицы кодов.

Максимальная скорость сжатия достигает порядка $10^8$ байт в секунду. Скорость разжатия зависит от структуры данных сильнее: для случайных данных она составляет порядка $10^7$ байт в секунду, а для периодических последовательностей может достигать $1.75 \cdot 10^8$ байт в секунду.

Установлено, что рост энтропии данных приводит к почти линейному ухудшению коэффициента сжатия, а также к снижению скорости архивирования и разархивирования. Полученные зависимости представлены графиками, построенными по результатам экспериментов.

Таблица с данными экспериментов:
\url{https://github.com/GidRadium/huffman-archiver/blob/main/research/benchmark_results.csv}.

Полученные графики:
\url{https://github.com/GidRadium/huffman-archiver/blob/main/research/plots}.

\section{Заключение}

В ходе выполнения работы получены следующие результаты:

\begin{itemize}
    \item выполнен обзор алгоритма Хаффмана и существующих подходов к построению архиваторов;
    \item разработано консольное приложение для сжатия и разжатия файлов по алгоритму Хаффмана;
    \item реализованы два режима архивирования, ориентированные на различные сценарии использования;
    \item разработан набор модульных и интеграционных тестов;
    \item проведено экспериментальное исследование эффективности и производительности реализации;
    \item подготовлена инфраструктура автоматической сборки, тестирования и генерации документации.
\end{itemize}

Репозиторий:
\url{https://github.com/GidRadium/huffman-archiver}.

Техническая документация:
\url{https://gidradium.github.io/huffman-archiver/}.

\end{document}
